\chapter{2023年全国甲卷（文）}

\section{单选题}

\begin{problem}
设全集 \(U=\{1,2,3,4,5\}\)，集合 \(M=\{1,4\}\)，\(N=\{2,5\}\)，则 \(N\cup\complement_U M=\)（\quad）

\choices
    {\(\{2,3,5\}\)}
    {\(\{1,3,4\}\)}
    {\(\{1,2,4,5\}\)}
    {\(\{2,3,4,5\}\)}

\end{problem}

\begin{answer}
A
\end{answer}

\begin{solution}
因为全集 \(U=\{1,2,3,4,5\}\)，集合 \(M=\{1,4\}\)，所以 \(\complement_U M=\{2,3,5\}\).
又 \(N=\{2,5\}\)，所以
\(N\cup\complement_U M=\{2,3,5\}\).
故选：\(A\).
\end{solution}

\begin{problem}
\(\dfrac{5\left(1+\i^3\right)}{\left(2+\i\right)\left(2-\i\right)}=\)（\quad）

\choices
    {\(-1\)}
    {\(1\)}
    {\(1-\i\)}
    {\(1+\i\)}

\end{problem}

\begin{answer}
C
\end{answer}

\begin{solution}
\(\dfrac{5\left(1+\i^3\right)}{\left(2+\i\right)\left(2-\i\right)} =\dfrac{5\left(1-\i\right)}{5} =1-\i\).
故选：\(C\).
\end{solution}

\begin{problem}
已知向量 \(\bs{a}=\left(3,1\right)\)，\(\bs{b}=\left(2,2\right)\)，则 \(\cos\left(\bs{a}+\bs{b},\bs{a}-\bs{b}\right)=\)（\quad）

\choices
    {\(\dfrac{1}{17}\)}
    {\(\dfrac{\sqrt{17}}{17}\)}
    {\(\dfrac{\sqrt{5}}{5}\)}
    {\(\dfrac{2\sqrt{5}}{5}\)}

\end{problem}

\begin{answer}
B
\end{answer}

\begin{solution}
因为 \(\bs{a}=\left(3,1\right)\)，\(\bs{b}=\left(2,2\right)\)，所以 \(\bs{a}+\bs{b}=\left(5,3\right)\)，\(\bs{a}-\bs{b}=\left(1,-1\right)\).
则
\(\left|\bs{a}+\bs{b}\right|^2=5^2+3^2=34\)，\(\left|\bs{a}-\bs{b}\right|^2=1^2+\left(-1\right)^2=2\)，
\(\left(\bs{a}+\bs{b}\right)\cdot\left(\bs{a}-\bs{b}\right)=5\times1+3\times\left(-1\right)=2\).
所以
\[
\cos\left(\bs{a}+\bs{b},\bs{a}-\bs{b}\right)=\dfrac{\left(\bs{a}+\bs{b}\right)\cdot\left(\bs{a}-\bs{b}\right)}{\left|\bs{a}+\bs{b}\right|\left|\bs{a}-\bs{b}\right|}=\dfrac{2}{\sqrt{34}\times\sqrt{2}}=\dfrac{\sqrt{17}}{17}
\]
故选：\(B\).
\end{solution}

\begin{problem}
某校文艺部有 4 名学生，其中高一，高二年级各 2 名. 从这 4 名学生中随机选 2 名组织校文艺汇演，则这 2 名学生来自不同年级的概率为（\quad）

\choices
    {\(\dfrac{1}{6}\)}
    {\(\dfrac{1}{3}\)}
    {\(\dfrac{1}{2}\)}
    {\(\dfrac{2}{3}\)}

\end{problem}

\begin{answer}
D
\end{answer}

\begin{solution}
依题意，从这 4 名学生中随机选 2 名组织校文艺汇演，总的基本事件有 \(C_4^2=6\) 件，
其中这 2 名学生来自不同年级的基本事件有 \(C_2^1C_2^1=4\) 件，
所以这 2 名学生来自不同年级的概率为
\(\dfrac{4}{6}=\dfrac{2}{3}\).
故选：\(D\).
\end{solution}

\begin{problem}
记 \(S_n\) 为等差数列 \(\{a_n\}\) 的前 \(n\) 项和. 若 \(a_2+a_6=10\)，\(a_4a_8=45\)，则 \(S_5=\)（\quad）

\choices
    {25}
    {22}
    {20}
    {15}

\end{problem}

\begin{answer}
C
\end{answer}

\begin{solution}
【法一】设等差数列 \(\{a_n\}\) 的公差为 \(d\)，首项为 \(a_1\)，依题意可得，
\(a_2+a_6=a_1+d+a_1+5d=10\)，
即
\(a_1+3d=5\).
又
\(a_4a_8=\left(a_1+3d\right)\left(a_1+7d\right)=45\)，
解得：\(d=1\)，\(a_1=2\)，
所以 \(S_5=5a_1+\dfrac{5\times4}{2}\times d=5\times2+10=20\).
故选：\(C\).

【法二】\(a_2+a_6=2a_4=10\)，\(a_4a_8=45\)，所以 \(a_4=5\)，\(a_8=9\)，
\(d=\dfrac{a_8-a_4}{8-4}=1\)，
于是 \(a_3=a_4-d=5-1=4\)，
所以
\(S_5=5a_3=20\).
故选：\(C\).
\end{solution}

\begin{problem}
执行下边的程序框图，则输出的 \(B=\)（\quad）

\choices
    {21}
    {34}
    {55}
    {89}

\end{problem}

\begin{answer}
B
\end{answer}

\begin{solution}
当 \(k=1\) 时，判断框条件满足，第一次执行循环体，\(A=1+2=3\)，\(B=3+2=5\)，\(k=1+1=2\)；
当 \(k=2\) 时，判断框条件满足，第二次执行循环体，\(A=3+5=8\)，\(B=8+5=13\)，\(k=2+1=3\)；
当 \(k=3\) 时，判断框条件满足，第三次执行循环体，\(A=8+13=21\)，\(B=21+13=34\)，\(k=3+1=4\)；
当 \(k=4\) 时，判断框条件不满足，跳出循环体，输出 \(B=34\).
故选：\(B\).
\end{solution}

\begin{problem}
设 \(F_1,F_2\) 为椭圆 \(C:\dfrac{x^2}{5}+y^2=1\) 的两个焦点，点 \(P\) 在 \(C\) 上，若 \(\overrightarrow{PF_1}\cdot\overrightarrow{PF_2}=0\)，则 \(\left|PF_1\right|\cdot\left|PF_2\right|=\)（\quad）

\choices
    {1}
    {2}
    {4}
    {5}

\end{problem}

\begin{answer}
B
\end{answer}

\begin{solution}
【法一】因为 \(\overrightarrow{PF_1}\cdot\overrightarrow{PF_2}=0\)，所以 \(\angle F_1PF_2=90^\circ\)，
从而 \(S_{\triangle F_1PF_2}=b^2\tan45^\circ=1=\dfrac{1}{2}\left|PF_1\right|\cdot\left|PF_2\right|\).
所以 \(\left|PF_1\right|\cdot\left|PF_2\right|=2\).
故选：\(B\).

【法二】因为 \(\overrightarrow{PF_1}\cdot\overrightarrow{PF_2}=0\)，所以 \(\angle F_1PF_2=90^\circ\)，由椭圆方程可知，
\(c^2=5-1=4 \Rightarrow c=2\)，
所以
\(\left|PF_1\right|^2+\left|PF_2\right|^2=\left|F_1F_2\right|^2=4^2=16\)，
又 \(\left|PF_1\right|+\left|PF_2\right|=2a=2\sqrt{5}\)，平方得：
\(\left|PF_1\right|^2+\left|PF_2\right|^2+2\left|PF_1\right|\left|PF_2\right|=20\)，
所以 \(\left|PF_1\right|\cdot\left|PF_2\right|=2\).
故选：\(B\).
\end{solution}

\begin{problem}
曲线 \(y=\dfrac{\e^x}{x+1}\) 在点 \(\left(1,\dfrac{\e}{2}\right)\) 处的切线方程为（\quad）

\choices
    {\(y=\dfrac{\e}{4}x\)}
    {\(y=\dfrac{\e}{2}x\)}
    {\(y=\dfrac{\e}{4}x+\dfrac{\e}{4}\)}
    {\(y=\dfrac{\e}{2}x+\dfrac{3\e}{4}\)}

\end{problem}

\begin{answer}
C
\end{answer}

\begin{solution}
设曲线 \(y=\dfrac{\e^x}{x+1}\) 在点 \(\left(1,\dfrac{\e}{2}\right)\) 处的切线方程为
\(y-\dfrac{\e}{2}=k\left(x-1\right)\).
因为
\(y=\dfrac{\e^x}{x+1}\)，
所以
\(y'=\dfrac{\e^x\left(x+1\right)-\e^x}{\left(x+1\right)^2} =\dfrac{x\e^x}{\left(x+1\right)^2}\)，
所以
\(k=y'\big|_{x=1}=\dfrac{\e}{4}\).
所以
\(y-\dfrac{\e}{2}=\dfrac{\e}{4}\left(x-1\right)\).
所以曲线 \(y=\dfrac{\e^x}{x+1}\) 在点 \(\left(1,\dfrac{\e}{2}\right)\) 处的切线方程为
\(y=\dfrac{\e}{4}x+\dfrac{\e}{4}\).
故选：\(C\).
\end{solution}

\begin{problem}
已知双曲线 \(\dfrac{x^2}{a^2}-\dfrac{y^2}{b^2}=1\left(a>0,b>0\right)\) 的离心率为 \(\sqrt{5}\)，其中一条渐近线与圆 \(\left(x-2\right)^2+\left(y-3\right)^2=1\) 交于 \(A,B\) 两点，则 \(|AB|=\)（\quad）

\choices
    {\(\dfrac{\sqrt{5}}{5}\)}
    {\(\dfrac{2\sqrt{5}}{5}\)}
    {\(\dfrac{3\sqrt{5}}{5}\)}
    {\(\dfrac{4\sqrt{5}}{5}\)}

\end{problem}

\begin{answer}
D
\end{answer}

\begin{solution}
由 \(e=\sqrt{5}\)，则
\(\dfrac{c^2}{a^2}=\dfrac{a^2+b^2}{a^2}=1+\dfrac{b^2}{a^2}=5\)，
解得
\(\dfrac{b}{a}=2\)，
所以双曲线的一条渐近线不妨取 \(y=2x\)，
则圆心 \((2,3)\) 到渐近线的距离
\(d=\dfrac{|2\times2-3|}{\sqrt{2^2+1}}=\dfrac{\sqrt{5}}{5}\)，
所以弦长
\(|AB|=2\sqrt{r^2-d^2}=2\sqrt{1-\dfrac{1}{5}}=\dfrac{4\sqrt{5}}{5}\).
故选：\(D\).
\end{solution}

\begin{problem}
在三棱锥 \(P-ABC\) 中，\(\triangle ABC\) 是边长为 \(2\) 的等边三角形，\(PA=PB=2\)，\(PC=\sqrt{6}\)，则该棱锥的体积为（\quad）

\bitmapfigure[max width=6.0cm,max totalheight=4.8cm]{img/2023/national_paper_a_liberal/q10_fig1.png}

\choices
    {1}
    {\(\sqrt{3}\)}
    {2}
    {3}

\end{problem}

\begin{answer}
A
\end{answer}

\begin{solution}
取 \(AB\) 中点 \(E\)，连接 \(PE\)，\(CE\)，如图，
因为 \(\triangle ABC\) 是边长为 \(2\) 的等边三角形，\(PA=PB=2\)，
所以 \(PE\perp AB\)，\(CE\perp AB\)，又 \(PE,CE\subset\) 平面 \(PEC\)，\(PE\cap CE=E\)，
所以 \(AB\perp\) 平面 \(PEC\)，
又
\(PE=CE=2\times\dfrac{\sqrt{3}}{2}=\sqrt{3}\)，\(PC=\sqrt{6}\)，
故
\(PC^2=PE^2+CE^2\)，
即 \(PE\perp CE\)，
所以
\[
V=V_{B-PEC}+V_{A-PEC}=\dfrac{1}{3}S_{\triangle PEC}\cdot AB=\dfrac{1}{3}\times\dfrac{1}{2}\times\sqrt{3}\times\sqrt{3}\times2=1
\]
故选：\(A\).
\end{solution}

\begin{problem}
已知函数 \(f(x)=\e^{-\left(x-1\right)^2}\). 记 \(a=f\left(\dfrac{\sqrt{2}}{2}\right)\)，\(b=f\left(\dfrac{\sqrt{3}}{2}\right)\)，\(c=f\left(\dfrac{\sqrt{6}}{2}\right)\)，则（\quad）

\choices
    {\(b>c>a\)}
    {\(b>a>c\)}
    {\(c>b>a\)}
    {\(c>a>b\)}

\end{problem}

\begin{answer}
A
\end{answer}

\begin{solution}
令 \(g(x)=-\left(x-1\right)^2\)，则 \(g(x)\) 开口向下，对称轴为 \(x=1\)，
因为
\(\dfrac{\sqrt{6}}{2}-1-\left(1-\dfrac{\sqrt{3}}{2}\right)=\dfrac{\sqrt{6}+\sqrt{3}}{2}-2\)，
而
\(\left(\sqrt{6}+\sqrt{3}\right)^2-4^2=9+6\sqrt{2}-16=6\sqrt{2}-7>0\)，
所以
\(\dfrac{\sqrt{6}}{2}-1-\left(1-\dfrac{\sqrt{3}}{2}\right)>0\)，
即
\(\dfrac{\sqrt{6}}{2}-1>1-\dfrac{\sqrt{3}}{2}\)，
所以由二次函数性质知
\(g\left(\dfrac{\sqrt{6}}{2}\right)<g\left(\dfrac{\sqrt{3}}{2}\right)\).
因为
\(\dfrac{\sqrt{6}}{2}-1-\left(1-\dfrac{\sqrt{2}}{2}\right)=\dfrac{\sqrt{6}+\sqrt{2}}{2}-2\)，
而
\(\left(\sqrt{6}+\sqrt{2}\right)^2-4^2=8+4\sqrt{3}-16=4\left(\sqrt{3}-2\right)<0\)，
即
\(\dfrac{\sqrt{6}}{2}-1<1-\dfrac{\sqrt{2}}{2}\)，
所以
\(g\left(\dfrac{\sqrt{6}}{2}\right)>g\left(\dfrac{\sqrt{2}}{2}\right)\).
综上，
\(g\left(\dfrac{\sqrt{2}}{2}\right)<g\left(\dfrac{\sqrt{6}}{2}\right)<g\left(\dfrac{\sqrt{3}}{2}\right)\).
又 \(y=\e^x\) 为增函数，故 \(a<c<b\)，即 \(b>c>a\).故选：\(A\).
\end{solution}

\begin{problem}
函数 \(y=f(x)\) 的图象由 \(y=\cos\left(2x+\dfrac{\pi}{6}\right)\) 的图象向左平移 \(\dfrac{\pi}{6}\) 个单位长度得到，则 \(y=f(x)\) 的图象与直线 \(y=\dfrac{1}{2}x-\dfrac{1}{2}\) 的交点个数为（\quad）

\choices
    {1}
    {2}
    {3}
    {4}

\end{problem}

\begin{answer}
C
\end{answer}

\begin{solution}
因为 \(y=\cos\left(2x+\dfrac{\pi}{6}\right)\) 向左平移 \(\dfrac{\pi}{6}\) 个单位所得函数为
\(y=\cos\left[2\left(x+\dfrac{\pi}{6}\right)+\dfrac{\pi}{6}\right]=\cos\left(2x+\dfrac{\pi}{2}\right)=-\sin2x\)，
所以 \(f(x)=-\sin2x\)，
当 \(x=0\) 时，\(y=\dfrac{1}{2}x-\dfrac{1}{2}=-\dfrac{1}{2}\)；当 \(x=1\) 时，\(y=\dfrac{1}{2}x-\dfrac{1}{2}=0\)，
所以直线 \(y=\dfrac{1}{2}x-\dfrac{1}{2}\) 过 \(\left(0,-\dfrac{1}{2}\right)\) 与 \((1,0)\) 两点，
考虑 \(2x=-\dfrac{3\pi}{2}\)，\(2x=\dfrac{3\pi}{2}\)，\(2x=\dfrac{7\pi}{2}\)，即 \(x=-\dfrac{3\pi}{4}\)，\(x=\dfrac{3\pi}{4}\)，\(x=\dfrac{7\pi}{4}\) 处 \(f(x)\) 与 \(y=\dfrac{1}{2}x-\dfrac{1}{2}\) 的大小关系，
当 \(x=-\dfrac{3\pi}{4}\) 时，
\(f\left(-\dfrac{3\pi}{4}\right)=-\sin\left(-\dfrac{3\pi}{2}\right)=-1\)，
\(y=\dfrac{1}{2}\times\left(-\dfrac{3\pi}{4}\right)-\dfrac{1}{2} =-\dfrac{3\pi+4}{8}<-1\)；
当 \(x=\dfrac{3\pi}{4}\) 时，
\(f\left(\dfrac{3\pi}{4}\right)=-\sin\dfrac{3\pi}{2}=1\)，
\(y=\dfrac{1}{2}\times\dfrac{3\pi}{4}-\dfrac{1}{2} =\dfrac{3\pi-4}{8}<1\)；
当 \(x=\dfrac{7\pi}{4}\) 时，
\(f\left(\dfrac{7\pi}{4}\right)=-\sin\dfrac{7\pi}{2}=1\)，
\(y=\dfrac{1}{2}\times\dfrac{7\pi}{4}-\dfrac{1}{2} =\dfrac{7\pi-4}{8}>1\).
所以 \(f(x)\) 与 \(y=\dfrac{1}{2}x-\dfrac{1}{2}\) 的交点个数为 \(3\).
故选：\(C\).
\end{solution}
\section{填空题}

\begin{problem}
记 \(S_n\) 为等比数列 \(\{a_n\}\) 的前 \(n\) 项和. 若 \(8S_6=7S_3\)，则 \(\{a_n\}\) 的公比为\fillinblank{}.

\end{problem}

\begin{answer}
\(-\dfrac{1}{2}\)
\end{answer}

\begin{solution}
若 \(q=1\)，
则由 \(8S_6=7S_3\) 得 \(8\cdot6a_1=7\cdot3a_1\)，则 \(a_1=0\)，不合题意.
所以 \(q\ne1\).
当 \(q\ne1\) 时，因为 \(8S_6=7S_3\)，
\(8\cdot\dfrac{a_1\left(1-q^6\right)}{1-q} =7\cdot\dfrac{a_1\left(1-q^3\right)}{1-q}\)，
即
\(8\left(1-q^6\right)=7\left(1-q^3\right)\)，
即
\(8\left(1+q^3\right)\left(1-q^3\right)=7\left(1-q^3\right)\)，
即
\(8\left(1+q^3\right)=7\)，
解得
\(q=-\dfrac{1}{2}\).
\end{solution}

\begin{problem}
若 \(f(x)=\left(x-1\right)^2+ax+\sin\left(x+\dfrac{\pi}{2}\right)\) 为偶函数，则 \(a=\fillinblank{}\).

\end{problem}

\begin{answer}
2
\end{answer}

\begin{solution}
因为 \(f(x)=\left(x-1\right)^2+ax+\sin\left(x+\dfrac{\pi}{2}\right)=\left(x-1\right)^2+ax+\cos x=x^2+\left(a-2\right)x+1+\cos x\)，
且函数为偶函数，
所以 \(a-2=0\)，解得 \(a=2\).
\end{solution}

\begin{problem}
若 \(x,y\) 满足约束条件 \(3x-2y\le3\)，\(-2x+3y\le3\)，\(x+y\ge1\)，则 \(z=3x+2y\) 的最大值为\fillinblank{}.

\end{problem}

\begin{answer}
15
\end{answer}

\begin{solution}
作出可行域，如图，

\bitmapfigure[max width=6cm,max totalheight=4.8cm]{img/2023/national_paper_a_liberal/q15_solution_fig1.png}

由图可知，当目标函数
\(y=-\dfrac{3}{2}x+\dfrac{z}{2}\)
过点 \(A\) 时，\(z\) 有最大值，
由 \(\begin{cases}
-2x+3y=3\\
3x-2y=3
\end{cases}\)
可得 \(\begin{cases}
x=3\\
y=3
\end{cases}\)
即 \(A(3,3)\)，
所以
\(z_{\max}=3\times3+2\times3=15\).
\end{solution}

\begin{problem}
在正方体 \(ABCD-A_1B_1C_1D_1\) 中，\(AB=4\)，\(O\) 为 \(AC_1\) 的中点，若该正方体的棱与球 \(O\) 的球面有公共点，则球 \(O\) 的半径的取值范围是\fillinblank{}.

\end{problem}

\begin{answer}
\([2\sqrt{2},2\sqrt{3}]\)
\end{answer}

\begin{solution}
设球的半径为 \(R\).
当球是正方体的外接球时，恰好经过正方体的每个顶点，所求的球的半径最大，若半径变得更大，球会包含正方体，导致球面和棱没有交点，
正方体的外接球直径 \(2R'\) 为体对角线长
\(AC_1=\sqrt{4^2+4^2+4^2}=4\sqrt{3}\)，
即 \(2R'=4\sqrt{3}\)，\(R'=2\sqrt{3}\)，故 \(R_{\max}=2\sqrt{3}\)；
分别取侧棱 \(AA_1\)，\(BB_1\)，\(CC_1\)，\(DD_1\) 的中点 \(M,H,G,N\).
设正方体的顶点坐标分别为 \(A(0,0,0)\)，\(B(4,0,0)\)，\(C(4,4,0)\)，\(D(0,4,0)\)，\(A_1(0,0,4)\)，\(B_1(4,0,4)\)，\(C_1(4,4,4)\)，\(D_1(0,4,4)\)，
则 \(M(0,0,2)\)，\(H(4,0,2)\)，\(G(4,4,2)\)，\(N(0,4,2)\)，所以四边形 \(MNGH\) 是边长为 \(4\) 的正方形.
又 \(O\) 为 \(AC_1\) 的中点，所以 \(O(2,2,2)\)，即 \(O\) 为正方形 \(MNGH\) 的对角线交点，
连接 \(MG\)，则
\(MG=4\sqrt{2}\)，
当球的一个大圆恰好是四边形 \(MNGH\) 的外接圆，球的半径达到最小，即 \(R\) 的最小值为 \(2\sqrt{2}\).
综上，
\(R\in[2\sqrt{2},2\sqrt{3}]\).
\end{solution}
\section{解答题}

\begin{problem}
记 \(\triangle ABC\) 的内角 \(A,B,C\) 的对边分别为 \(a,b,c\)，已知 \(\dfrac{b^2+c^2-a^2}{\cos A}=2\).
\begin{enumerate}
\item 求 \(bc\)；
\item 若 \(\dfrac{a\cos B-b\cos A}{a\cos B+b\cos A}-\dfrac{b}{c}=1\)，求 \(\triangle ABC\) 面积.
\end{enumerate}

\end{problem}

\begin{solution}
（1）因为
\(a^2=b^2+c^2-2bc\cos A\)，
所以
\(\dfrac{b^2+c^2-a^2}{\cos A} =\dfrac{2bc\cos A}{\cos A} =2bc=2\)，解得 \(bc=1\).

（2）由正弦定理可得
\(\dfrac{a\cos B-b\cos A}{a\cos B+b\cos A}-\dfrac{b}{c}
=\dfrac{\sin A\cos B-\sin B\cos A}{\sin A\cos B+\sin B\cos A}-\dfrac{\sin B}{\sin C}
=\dfrac{\sin\left(A-B\right)}{\sin\left(A+B\right)}-\dfrac{\sin B}{\sin\left(A+B\right)}
=\dfrac{\sin\left(A-B\right)-\sin B}{\sin\left(A+B\right)}
=1\)，
变形可得：
\(\sin\left(A-B\right)-\sin\left(A+B\right)=\sin B\)，
即
\(-2\cos A\sin B=\sin B\)，
而 \(0<\sin B\le1\)，所以
\(\cos A=-\dfrac{1}{2}\)，
又 \(0<A<\pi\)，所以
\(\sin A=\dfrac{\sqrt{3}}{2}\)，
故 \(\triangle ABC\) 的面积为 \(S_{\triangle ABC}=\dfrac{1}{2}bc\sin A=\dfrac{1}{2}\times1\times\dfrac{\sqrt{3}}{2}=\dfrac{\sqrt{3}}{4}\).
\end{solution}

\begin{problem}
如图，在三棱柱 \(ABC-A_1B_1C_1\) 中，\(A_1C\perp\) 平面 \(ABC\)，\(\angle ACB=90^\circ\).

\begin{enumerate}
\item 证明：平面 \(ACC_1A_1\perp\) 平面 \(BCC_1B_1\)；
\item 设 \(AB=A_1B\)，\(AA_1=2\)，求四棱锥 \(A_1-BB_1C_1C\) 的高.
\end{enumerate}

\bitmapfigure[max width=6.0cm,max totalheight=4.8cm]{img/2023/national_paper_a_liberal/q18_fig1.png}

\end{problem}

\begin{solution}
（1）证明：因为 \(A_1C\perp\) 平面 \(ABC\)，\(BC\subset\) 平面 \(ABC\)，
所以 \(A_1C\perp BC\)，
又因为 \(\angle ACB=90^\circ\)，即 \(AC\perp BC\)，
\(A_1C,AC\subset\) 平面 \(ACC_1A_1\)，\(A_1C\cap AC=C\)，
所以 \(BC\perp\) 平面 \(ACC_1A_1\)，
又因为 \(BC\subset\) 平面 \(BCC_1B_1\)，
所以平面 \(ACC_1A_1\perp\) 平面 \(BCC_1B_1\).

（2）如图，

\bitmapfigure[max width=3.8cm,max totalheight=4.8cm]{img/2023/national_paper_a_liberal/q18_solution_fig1.png}

过点 \(A_1\) 作 \(A_1O\perp CC_1\)，垂足为 \(O\).
因为平面 \(ACC_1A_1\perp\) 平面 \(BCC_1B_1\)，平面 \(ACC_1A_1\cap\) 平面 \(BCC_1B_1=CC_1\)，\(A_1O\subset\) 平面 \(ACC_1A_1\)，
所以 \(A_1O\perp\) 平面 \(BCC_1B_1\)，
所以四棱锥 \(A_1-BB_1C_1C\) 的高为 \(A_1O\).
因为 \(A_1C\perp\) 平面 \(ABC\)，\(AC,BC\subset\) 平面 \(ABC\)，
所以 \(A_1C\perp BC\)，\(A_1C\perp AC\)，
又因为 \(A_1B=AB\)，\(BC\) 为公共边，
所以 \(\triangle ABC\) 与 \(\triangle A_1BC\) 全等，所以 \(A_1C=AC\).
设 \(A_1C=AC=x\)，则 \(A_1C_1=x\)，
所以 \(O\) 为 \(CC_1\) 中点，
\(OC_1=\dfrac{1}{2}AA_1=1\)，
又因为 \(A_1C\perp AC\)，所以
\(A_1C^2+AC^2=AA_1^2\)，
即
\(x^2+x^2=2^2\)，
解得
\(x=\sqrt{2}\)，
所以 \(A_1O=\sqrt{A_1C_1^2-OC_1^2}=\sqrt{\left(\sqrt{2}\right)^2-1^2}=1\)，
所以四棱锥 \(A_1-BB_1C_1C\) 的高为 \(1\).
\end{solution}

\begin{problem}
一项试验旨在研究臭氧效应，试验方案如下：选 40 只小白鼠，随机地将其中 20 只分配到试验组，另外 20 只分配到对照组，试验组的小白鼠饲养在高浓度臭氧环境，对照组的小白鼠饲养在正常环境，一段时间后统计每只小白鼠体重的增加量（单位：g）. 试验结果如下：

对照组的小白鼠体重的增加量从小到大排序为

15.2，18.8，20.2，21.3，22.5，23.2，25.8，26.5，27.5，30.1，

32.6，34.3，34.8，35.6，35.6，35.8，36.2，37.3，40.5，43.2

试验组的小白鼠体重的增加量从小到大排序为

7.8，9.2，11.4，12.4，13.2，15.5，16.5，18.0，18.8，19.2，

19.8，20.2，21.6，22.8，23.6，23.9，25.1，28.2，32.3，36.5
\begin{enumerate}
\item 计算试验组的样本平均数；
\item 求 40 只小白鼠体重的增加量的中位数 \(m\)，再分别统计两样本中小于 \(m\) 与不小于 \(m\) 的数据的个数，完成如下列联表
\end{enumerate}

\begin{center}
\begin{tabular}{|c|c|c|}
\hline
& \(<m\) & \(\ge m\) \\
\hline
对照组 &  &  \\
\hline
试验组 &  &  \\
\hline
\end{tabular}
\end{center}

\begin{enumerate}
\setcounter{enumi}{2}
\item 根据（2）中的列联表，能否有 \(95\%\) 的把握认为小白鼠在高浓度臭氧环境中与在正常环境中体重的增加量有差异？
\end{enumerate}

附：\(K^2=\dfrac{n\left(ad-bc\right)^2}{\left(a+b\right)\left(c+d\right)\left(a+c\right)\left(b+d\right)}\)，

\begin{center}
\begin{tabular}{c|ccc}
\hline
\(P\left(K^2\ge k\right)\) & 0.100 & 0.050 & 0.010 \\
\hline
\(k\) & 2.706 & 3.841 & 6.635 \\
\hline
\end{tabular}
\end{center}

\end{problem}

\begin{solution}
（1）试验组样本平均数为：
\[
\begin{aligned}
\dfrac{1}{20}\bigl(&7.8+9.2+11.4+12.4+13.2+15.5+16.5\\
&+18.0+18.8+19.2+19.8+20.2+21.6+22.8\\
&+23.6+23.9+25.1+28.2+32.3+36.5\bigr)\\
&=\dfrac{396}{20}=19.8.
\end{aligned}
\]

（2）依题意，可知这 40 只小鼠体重的中位数是将两组数据合在一起，从小到大排后第 20 位与第 21 位数据的平均数，
由原数据可得第 11 位数据为 18.8，后续依次为 19.2，19.8，20.2，20.2，21.3，21.6，22.5，22.8，23.2，23.6，\(\cdots\)，
故第 20 位为 23.2，第 21 位数据为 23.6，
所以
\(m=\dfrac{23.2+23.6}{2}=23.4\)，
故列联表为：

\begin{center}
\begin{tabular}{c|ccc}
\hline
& \(<m\) & \(\ge m\) & 合计 \\
\hline
对照组 & 6 & 14 & 20 \\
试验组 & 14 & 6 & 20 \\
\hline
合计 & 20 & 20 & 40 \\
\hline
\end{tabular}
\end{center}

（3）由（2）可得，
\(K^2=\dfrac{40\times\left(6\times6-14\times14\right)^2}{20\times20\times20\times20}=6.400>3.841\)，
所以能有 \(95\%\) 的把握认为小白鼠在高浓度臭氧环境中与在正常环境中体重的增加量有差异.
\end{solution}

\begin{problem}
已知函数 \(f(x)=ax-\dfrac{\sin x}{\cos^2x}\)，\(x\in\left(0,\dfrac{\pi}{2}\right)\).
\begin{enumerate}
\item 当 \(a=1\) 时，讨论 \(f(x)\) 的单调性；
\item 若 \(f(x)+\sin x<0\)，求 \(a\) 的取值范围.
\end{enumerate}

\end{problem}

\begin{solution}
（1）因为 \(a=1\)，所以
\(f(x)=x-\dfrac{\sin x}{\cos^2x}\)，\(x\in\left(0,\dfrac{\pi}{2}\right)\)，
则
\begin{align*}
f'(x)&=1-\dfrac{\cos x\cos^2x-2\cos x\left(-\sin x\right)\sin x}{\cos^4x}\\
&=1-\dfrac{\cos^2x+2\sin^2x}{\cos^3x}\\
&=\dfrac{\cos^3x-\cos^2x-2\left(1-\cos^2x\right)}{\cos^3x}\\
&=\dfrac{\cos^3x+\cos^2x-2}{\cos^3x}
\end{align*}
令 \(t=\cos x\)，由于 \(x\in\left(0,\dfrac{\pi}{2}\right)\)，所以 \(t=\cos x\in(0,1)\)，
所以
\[
\cos^3x+\cos^2x-2=t^3+t^2-2=t^3-t^2+2t^2-2=t^2\left(t-1\right)+2\left(t+1\right)\left(t-1\right)=\left(t^2+2t+2\right)\left(t-1\right)
\]
因为
\(t^2+2t+2=\left(t+1\right)^2+1>0\)，\(t-1<0\)，\(\cos^3x=t^3>0\)，
所以
\(f'(x)=\dfrac{\cos^3x+\cos^2x-2}{\cos^3x}<0\)
在 \(\left(0,\dfrac{\pi}{2}\right)\) 上恒成立，
所以 \(f(x)\) 在 \(\left(0,\dfrac{\pi}{2}\right)\) 上单调递减.

（2）法一：
构建
\(g(x)=f(x)+\sin x=ax-\dfrac{\sin x}{\cos^2x}+\sin x\)（\(0<x<\dfrac{\pi}{2}\)），
则
\(g'(x)=a-\dfrac{1+\sin^2x}{\cos^3x}+\cos x\)（\(0<x<\dfrac{\pi}{2}\)），
若 \(g(x)=f(x)+\sin x<0\)，且 \(g(0)=f(0)+\sin0=0\)，
则
\(g'(0)=a-1+1=a\le0\)，
解得 \(a\le0\)，
当 \(a=0\) 时，因为
\(\sin x-\dfrac{\sin x}{\cos^2x} =\sin x\left(1-\dfrac{1}{\cos^2x}\right)\)，
又 \(x\in\left(0,\dfrac{\pi}{2}\right)\)，所以 \(0<\sin x<1\)，\(0<\cos x<1\)，则 \(\dfrac{1}{\cos^2x}>1\)，
所以
\(f(x)+\sin x=\sin x-\dfrac{\sin x}{\cos^2x}<0\)，
满足题意；
当 \(a<0\) 时，由于 \(0<x<\dfrac{\pi}{2}\)，所以 \(x>0\)，从而 \(ax<0\)，
所以 \(f(x)+\sin x=ax-\dfrac{\sin x}{\cos^2x}+\sin x<\sin x-\dfrac{\sin x}{\cos^2x}<0\).
满足题意；
综上所述：若 \(f(x)+\sin x<0\)，等价于 \(a\le0\)，
所以 \(a\) 的取值范围为 \((-\infty,0]\).

法二：
因为
\[
\sin x-\dfrac{\sin x}{\cos^2x}=\dfrac{\sin x\cos^2x-\sin x}{\cos^2x}=\dfrac{\sin x\left(\cos^2x-1\right)}{\cos^2x}=-\dfrac{\sin^3x}{\cos^2x}
\]
因为 \(x\in\left(0,\dfrac{\pi}{2}\right)\)，所以 \(0<\sin x<1\)，\(0<\cos x<1\)，
故
\(\sin x-\dfrac{\sin x}{\cos^2x}<0\)
在 \(\left(0,\dfrac{\pi}{2}\right)\) 上恒成立，
所以当 \(a=0\) 时，
\(f(x)+\sin x=\sin x-\dfrac{\sin x}{\cos^2x}<0\)，
满足题意；
当 \(a<0\) 时，由于 \(0<x<\dfrac{\pi}{2}\)，所以 \(x>0\)，从而 \(ax<0\)，
所以 \(f(x)+\sin x=ax-\dfrac{\sin x}{\cos^2x}+\sin x<\sin x-\dfrac{\sin x}{\cos^2x}<0\).
满足题意；
当 \(a>0\) 时，因为
\(f(x)+\sin x=ax-\dfrac{\sin x}{\cos^2x}+\sin x =ax-\dfrac{\sin^3x}{\cos^2x}\)，
令
\(g(x)=ax-\dfrac{\sin^3x}{\cos^2x}\)（\(0<x<\dfrac{\pi}{2}\)），
则
\(g'(x)=a-\dfrac{3\sin^2x\cos^2x+2\sin^4x}{\cos^3x}\)，
注意到
\(g'(0)=a-\dfrac{3\sin^20\cos^20+2\sin^40}{\cos^30}=a>0\)，
若 \(\forall0<x<\dfrac{\pi}{2}\)，\(g'(x)>0\)，则 \(g(x)\) 在 \(\left(0,\dfrac{\pi}{2}\right)\) 上单调递增，
注意到 \(g(0)=0\)，所以 \(g(x)>g(0)=0\)，即 \(f(x)+\sin x>0\)，不满足题意；
若 \(\exists0<x_0<\dfrac{\pi}{2}\)，\(g'(x_0)<0\)，则 \(g'(0)g'(x_0)<0\)，
所以在 \(\left(0,\dfrac{\pi}{2}\right)\) 上最靠近 \(x=0\) 处必存在零点 \(x_1\in\left(0,\dfrac{\pi}{2}\right)\)，使得 \(g'(x_1)=0\)，
此时 \(g'(x)\) 在 \((0,x_1)\) 上有 \(g'(x)>0\)，所以 \(g(x)\) 在 \((0,x_1)\) 上单调递增，
则在 \((0,x_1)\) 上有 \(g(x)>g(0)=0\)，即 \(f(x)+\sin x>0\)，不满足题意；
综上：\(a\le0\).
\end{solution}

\begin{problem}
已知直线 \(x-2y+1=0\) 与抛物线 \(C:y^2=2px\left(p>0\right)\) 交于 \(A,B\) 两点，\(|AB|=4\sqrt{15}\).
\begin{enumerate}
\item 求 \(p\)；
\item 设 \(F\) 为 \(C\) 的焦点，\(M,N\) 为 \(C\) 上两点，且 \(\overrightarrow{FM}\cdot\overrightarrow{FN}=0\)，求 \(\triangle MFN\) 面积的最小值.
\end{enumerate}

\end{problem}

\begin{solution}
（1）设 \(A\left(x_A,y_A\right)\)，\(B\left(x_B,y_B\right)\)，
由 \(\begin{cases}
x-2y+1=0\\
y^2=2px
\end{cases}\)
可得，
\(y^2-4py+2p=0\)，
所以
\(y_A+y_B=4p\)，\(y_Ay_B=2p\)，
所以
\[
|AB|=\sqrt{\left(x_A-x_B\right)^2+\left(y_A-y_B\right)^2}=\sqrt{5}\left|y_A-y_B\right|=\sqrt{5}\sqrt{\left(y_A+y_B\right)^2-4y_Ay_B}=4\sqrt{15}
\]
即
\(2p^2-p-6=0\)，
因为 \(p>0\)，解得：
\(p=2\).

（2）因为 \(F(1,0)\)，若直线 \(MN\) 的斜率为零，则 \(M,N\) 的纵坐标相等.
又 \(M,N\) 都在抛物线 \(y^2=4x\) 上，所以 \(M,N\) 的横坐标也相等，这与 \(M,N\) 为两点矛盾，
所以直线 \(MN\) 的斜率不可能为零，
设直线 \(MN:x=my+n\)，\(M\left(x_1,y_1\right)\)，\(N\left(x_2,y_2\right)\)，
由 \(\begin{cases}
y^2=4x\\
x=my+n
\end{cases}\)
可得，
\(y^2-4my-4n=0\)，
所以，
\(y_1+y_2=4m\)，\(y_1y_2=-4n\)，
\(\Delta=16m^2+16n>0 \Rightarrow m^2+n>0\)，
因为 \(\overrightarrow{FM}\cdot\overrightarrow{FN}=0\)，所以
\(\left(x_1-1\right)\left(x_2-1\right)+y_1y_2=0\)，
即
\(\left(my_1+n-1\right)\left(my_2+n-1\right)+y_1y_2=0\)，
亦即
\(\left(m^2+1\right)y_1y_2+m\left(n-1\right)\left(y_1+y_2\right)+\left(n-1\right)^2=0\)，
将 \(y_1+y_2=4m\)，\(y_1y_2=-4n\) 代入得，
\(4m^2=n^2-6n+1\)，\(4\left(m^2+n\right)=\left(n-1\right)^2>0\)，
所以 \(n\ne1\)，且
\(n^2-6n+1\ge0\)，
解得
\(n\ge3+2\sqrt{2}\text{ 或 }n\le3-2\sqrt{2}\).
设点 \(F\) 到直线 \(MN\) 的距离为 \(d\)，所以
\(d=\dfrac{|n-1|}{\sqrt{1+m^2}}\)，
\(|MN| =\sqrt{\left(x_1-x_2\right)^2+\left(y_1-y_2\right)^2} =\sqrt{1+m^2}\left|y_1-y_2\right|\)
\(=\sqrt{1+m^2}\sqrt{16m^2+16n} =2\sqrt{1+m^2}\sqrt{4\left(n^2-6n+1\right)+16n}\)
\(=2\sqrt{1+m^2}|n-1|\)，
所以 \(\triangle MNF\) 的面积 \(S=\dfrac{1}{2}\times|MN|\times d=\dfrac{1}{2}\times\dfrac{|n-1|}{\sqrt{1+m^2}}\times2\sqrt{1+m^2}|n-1|=\left(n-1\right)^2\)，
而 \(n\ge3+2\sqrt{2}\) 或 \(n\le3-2\sqrt{2}\)，所以，
\(\text{当 }n=3-2\sqrt{2}\text{ 时，}\triangle MNF\text{ 的面积 }S_{\min} =\left(2-2\sqrt{2}\right)^2 =12-8\sqrt{2}\).
\end{solution}

\begin{problem}
已知点 \(P(2,1)\)，直线 \(l\) 的参数方程为 \(x=2+t\cos\alpha\)，\(y=1+t\sin\alpha\)（\(t\) 为参数），\(\alpha\) 为 \(l\) 的倾斜角，\(l\) 与 \(x\) 轴正半轴，\(y\) 轴正半轴分别交于 \(A,B\)，且 \(|PA|\cdot|PB|=4\).
\begin{enumerate}
\item 求 \(\alpha\)；
\item 以坐标原点为极点，\(x\) 轴正半轴为极轴建立极坐标系，求 \(l\) 的极坐标方程.
\end{enumerate}

\end{problem}

\begin{solution}
（1）因为 \(l\) 与 \(x\) 轴，\(y\) 轴正半轴交于 \(A,B\) 两点，所以
\(\dfrac{\pi}{2}<\alpha<\pi\)，
令 \(x=0\)，
\(t_1=-\dfrac{2}{\cos\alpha}\)，
令 \(y=0\)，
\(t_2=-\dfrac{1}{\sin\alpha}\)，
所以
\(|PA||PB|=|t_2t_1|=\left|\dfrac{2}{\sin\alpha\cos\alpha}\right| =\left|\dfrac{4}{\sin2\alpha}\right|=4\)，
所以
\(\sin2\alpha=\pm1\)，
即
\(2\alpha=\dfrac{\pi}{2}+k\pi\)，
解得
\(\alpha=\dfrac{\pi}{4}+\dfrac{1}{2}k\pi\)，\(k\in \Z\)，
因为 \(\dfrac{\pi}{2}<\alpha<\pi\)，所以
\(\alpha=\dfrac{3\pi}{4}\).

（2）由（1）可知，直线 \(l\) 的斜率为 \(\tan\alpha=-1\)，且过点 \((2,1)\)，
所以直线 \(l\) 的普通方程为：
\(y-1=-\left(x-2\right)\)，
即
\(x+y-3=0\)，
由
\(x=\rho\cos\theta\)，\(y=\rho\sin\theta\)
可得直线 \(l\) 的极坐标方程为
\(\rho\cos\theta+\rho\sin\theta-3=0\).
\end{solution}

\begin{problem}
已知 \(f(x)=2|x-a|-a\)，\(a>0\).
\begin{enumerate}
\item 求不等式 \(f(x)<x\) 的解集；
\item 若曲线 \(y=f(x)\) 与 \(x\) 轴所围成的图形的面积为 \(2\)，求 \(a\).
\end{enumerate}

\end{problem}

\begin{solution}
（1）若 \(x\le a\)，则
\(f(x)=2a-2x-a<x\)，
即
\(3x>a\)，
解得
\(x>\dfrac{a}{3}\)，
即
\(\dfrac{a}{3}<x\le a\)，
若 \(x>a\)，则
\(f(x)=2x-2a-a<x\)，
解得
\(x<3a\)，
即
\(a<x<3a\)，
综上，不等式的解集为
\(\left(\dfrac{a}{3},3a\right)\).

（2）
\[
f(x)=\begin{cases}
-2x+a,&x\le a\\
2x-3a,&x>a
\end{cases}
\]
画出 \(f(x)\) 的草图，则 \(f(x)\) 与坐标轴围成 \(\triangle ADO\) 与 \(\triangle ABC\)，
\(\triangle ABC\) 的高为 \(a\)，\(D(0,a)\)，\(A\left(\dfrac{a}{2},0\right)\)，\(B\left(\dfrac{3a}{2},0\right)\)，所以 \(|AB|=a\)，
所以
\(S_{\triangle OAD}+S_{\triangle ABC} =\dfrac{1}{2}|OA|\cdot a+\dfrac{1}{2}|AB|\cdot a =\dfrac{3}{4}a^2=2\)，
解得
\(a=\dfrac{2\sqrt{6}}{3}\).
\end{solution}
